% 三种收敛方式的蕴含关系与反例（教学示意，无数据）。
% 实线：已证明的蕴含（有界收敛定理、Markov不等式、Jensen不等式，1<=p<q）。
% 红色虚线：反向蕴含不成立，旁注所用反例：
%   依概率 -/-> 几乎处处：二进区间扫描；几乎处处 -/-> p阶平均：尖峰 n I{U<1/n}；
%   p阶平均 -/-> q阶平均：n^{1/q} I{U<1/n}（p阶误差 n^{p/q-1}->0，q阶误差恒为1）。
\begin{tikzpicture}[line width=.6pt,font=\fontsize{8}{10}\selectfont,>=Stealth,
  mode/.style={draw=FigureBlue,fill=FigureBlueFill,rounded corners=1.8pt,minimum width=22mm,minimum height=7.5mm,inner sep=2pt,text=FigureInk},
  imp/.style={->,draw=FigureStroke,line width=.75pt},
  fail/.style={->,draw=FigureRed,dashed,line width=.7pt},
  lab/.style={font=\fontsize{7.5}{9}\selectfont,text=FigureMuted},
  flab/.style={font=\fontsize{7.5}{9}\selectfont,text=FigureRed}]
 \node[mode] (as) at (0,0) {几乎处处收敛};
 \node[mode] (pr) at (44mm,0) {依概率收敛};
 \node[mode] (lp) at (88mm,0) {$p$阶平均收敛};
 \node[mode] (lq) at (88mm,24mm) {$q$阶平均收敛};
 \draw[imp] (as) -- node[lab,below=1pt] {有界收敛定理} (pr);
 \draw[imp] (lp) -- node[lab,below=1pt] {Markov不等式} (pr);
 \draw[imp] (lq) -- node[lab,left=1pt] {Jensen不等式} (lp);
 \draw[fail] (pr.north) to[out=135,in=45] node[flab,above=1pt] {反例：二进区间扫描} (as.north);
 \draw[fail] (as.south) to[out=-40,in=-140] node[flab,below=1pt] {反例：尖峰$n\,\mathbb{I}_{\{U<1/n\}}$} (lp.south);
 \draw[fail] (lp.east) to[out=20,in=-20] node[flab,right=2pt,align=left] {反例：\\$n^{1/q}\,\mathbb{I}_{\{U<1/n\}}$} (lq.east);
 \node[lab,anchor=east] at ([xshift=-3mm]lq.west) {$1\leq p<q<\infty$};
\end{tikzpicture}
